% 中英文摘要章节
% 本科：中英文摘要排在摘要页内，中文在前、英文在后（本科指导手册「论文装订顺序」要求同页）；
% 硕士/博士：中文、英文摘要各自另起一页。
% 前置部分的罗马页码由 main.tex 的 \frontmatter 统一初始化并连续编号，此处不再重置，
% 否则中文摘要跨页时页码会倒退出现重复的 I（issue #7）。
\zihao{-4}
\vspace*{-10mm}

\begin{center}
  \heiti\zihao{3}\textmd{\ifcsdef{thesisTitleAbs}{\thesisTitleAbs}{\thesisTitle}}
  \vspace{2mm}
  \ifhaveSubTitle
  \begin{spacing}{1.2}
    \sihao\selectfont{\textmd{\kaishu{-----}\thesisSubTitle}}
  \end{spacing}
  \fi
\end{center}

\vspace*{0mm}

{\let\clearpage\relax \chapter*{\textmd{}}}
% 加入目录
\phantomsection
\addcontentsline{toc}{chapter}{摘~~~~要}

\vspace*{-12mm}
\setstretch{1.53}
\setlength{\parskip}{0em}

\textbf{\heiti\xiaosi{摘要：}}%
% 中文摘要正文从这里开始-----------------------------------------------------------------
{\kaishu{\expandafter\ignorespaces\abstractCN}}

\vspace{1em}
\textbf{\heiti\xiaosi{关键词：}}%
{\kaishu{\expandafter\ignorespaces\keywordsCN}}

% 英文摘要章节
% 硕士/博士：英文摘要另起一页，页首留出与中文摘要页一致的上边距；
% 本科：与中文摘要同页、英文在后（见文件开头说明）
\ifzufeBachelor
  \vspace*{15mm}
\else
  \newpage
  \vspace*{-2.5mm}
\fi

\begin{spacing}{0.95}
  \centering
  \zihao{3}\textbf{\thesisTitleEN}
  \vspace{2mm}
  \ifhaveSubTitle
  \begin{spacing}{1.2}
    \sihao\selectfont{\textmd{{-----}\thesisSubTitleEN}}
  \end{spacing}
  \fi
\end{spacing}

\vspace*{0mm}

{\let\clearpage\relax \chapter*{\zihao{-3}\textmd{}\vskip -3bp}}
\phantomsection
\addcontentsline{toc}{chapter}{Abstract}

\setstretch{1.53}
\setlength{\parskip}{0em}

% 四个标签都只保留一个词间空格：标签行末换行和 \abstractXX/\keywordsXX
% 宏体首行的换行都会各吐一个空格，用 \ignorespaces 吞掉，否则标签与正文之
% 间会排出 2~3 个空格宽的空隙（issue #10）。
% 中文侧同样需要：xeCJK 会吃掉 CJK 字符间的单个空格，仅靠换行时中文摘要看
% 不出问题；但正文若以拉丁字符开头（如「R&D 效率」），多余空格就会显形。
\textbf{\heiti\xiaosi{Abstract:}}\space
\expandafter\ignorespaces\abstractEN


\textbf{\xiaosi{Key words:}}\space
\expandafter\ignorespaces\keywordsEN
\newpage
